%---------------------------------------------------------------------------------
%                西南交通大学研究生学位论文：第二章内容
%---------------------------------------------------------------------------------
\chapter{相关理论与技术}

本章对与本文研究相关的理论基础和关键技术进行讲述，重点讨论大语言模型与工具增强、可靠查询生成、数据分析规划与结果生成，以及系统编排与服务协同等内容。
相关研究表明，面向政务问数与数据分析的智能系统已不再是单一文本生成问题，而是由语义理解、结构化生成、外部工具执行、流程组织和结果表达共同构成的复合技术链\cite{zhao2023llmsurvey,王昀2024}。
基于这一认识，本章按照模型基础，查询技术，分析技术，系统协同的思路展开，旨在为全文研究提供必要的知识背景和方法参照。

\section{大语言模型与工具增强基础}

大语言模型（Large Language Model，LLM）是在大规模语料上进行预训练，并通过指令对齐获得任务遵循能力的语言模型\cite{zhao2023llmsurvey}。
在面向数据查询、统计分析和服务编排的应用中，大语言模型的价值并不只体现在通用文本生成能力上，更体现在统一语义建模、少样本任务适配、结构化推理以及与外部工具协同等方面。
这些能力共同奠定了复杂智能系统的模型基础\cite{王昀2024,朱炫鹏2024}。

\subsection{Transformer 与指令对齐}

Transformer以自注意力机制为核心，通过并行建模序列中不同位置之间的依赖关系，使模型能够在统一表示空间内同时处理词语语义、上下文关系和结构约束\cite{vaswani2017Transformer_pages}。
这类架构具有较强的长距离依赖建模能力，因此更适合承载自然语言理解、结构化生成和多步骤推理等任务。

从技术演进看，GPT类解码器模型\cite{radford2018gpt,radford2019gpt2,brown2020gpt3,openai2023gpt4}强化了自回归生成能力，BERT类模型\cite{devlin2019bert}强化了双向语义表示能力，T5\cite{raffel2020t5}则提供了统一的文本到文本建模范式。
朱炫鹏等人\cite{朱炫鹏2024}指出，当前大语言模型的发展一方面沿着模型能力增强持续推进，另一方面也在通过推理效率优化、架构改进和工程化手段提升实际可用性。
对于自然语言到SQL、自动化分析和多步骤任务执行而言，解码器式或统一生成式模型通常更具现实适用性，因为其输出目标往往天然表现为SQL语句、分析计划、代码片段或结构化结果。

仅有预训练还不足以支撑复杂应用。
指令微调（Instruction Tuning）通过构造“指令 - 输入 - 输出”样本，使模型学会遵循任务边界、输出格式和领域约束\cite{wei2022flan,chung2022flanT5}；
进一步的对齐训练，如基于人类反馈的强化学习\cite{ouyang2022instructgpt}，则使模型输出更符合安全性、可用性和稳定性要求。
对于需要兼顾表达规范、统计口径和执行可靠性的应用场景而言，这种“预训练 + 任务对齐”的技术路线尤为重要\cite{王昀2024}。

% todo

\subsection{上下文学习、链式推理与结构化输出}

上下文学习（In-Context Learning，ICL）使模型能够在不更新参数的情况下，仅通过提示中的任务说明、数据库模式信息和少量示例完成新任务\cite{brown2020gpt3,dong2023icl_survey}。
这使模型能够在不同数据库、不同任务模板和不同分析需求之间快速迁移，而不必为每一种场景重新训练专用模型。
在复杂数据应用中，模式描述、历史样例和输出格式约束通常都会作为上下文的一部分被显式组织进提示中，从而提升模型对当前任务的适配性。

链式推理（Chain-of-Thought，CoT）进一步提升了复杂任务中的中间步骤显式化能力\cite{wei2022chain,kojima2022zeroshotcot}。
当用户需求涉及同比、占比、排名、条件筛选和多阶段聚合等隐含逻辑时，模型若直接输出最终SQL或最终结论，容易在中间语义转换中出现遗漏和幻觉。
通过显式展开推理步骤，模型可以先识别指标、维度和约束，再逐步构造查询或分析计划，从而提升复杂任务的稳定性。
自一致性\cite{wang2023selfconsistency}和思维树\cite{yao2023tot}等方法进一步说明，多路径推理与结果比较能够有效缓解单路径生成的不稳定问题。

结构化输出则把模型能力从“自由文本生成”进一步推进到“受约束结果生成”。
在数据系统中，输出对象通常不是普通自然语言段落，而是SQL语句、JSON计划、代码调用参数或固定槽位结果。
因此，结构化输出并不是简单的后处理步骤，而是大模型推理过程的一部分。
只有当模型能够在提示约束下生成字段完整、格式稳定、语义清晰的中间结果时，后续执行、校验和编排才具备可靠基础\cite{tang2024struc}。

\subsection{工具调用与代码执行}

大语言模型擅长语义理解与生成，但在精确计算、数据库访问、代码执行和外部状态感知方面存在明显局限\cite{qin2024toollearning}。
因此，当前实用系统普遍采用“模型负责推理，工具负责执行”的增强范式。
Toolformer\cite{schick2024toolformer}、Gorilla\cite{patil2023gorilla}等研究表明，模型能够学习何时调用外部工具以及如何构造调用参数；
ReAct\cite{yao2023react}进一步把推理、行动和观察组织成闭环，使模型能够根据执行结果持续修正自身决策。

这一范式对于数据类任务尤为关键。
在查询生成场景中，SQL需要交由数据库执行并根据报错信息或执行结果进行修正；
在数据分析场景中，模型生成的计划往往需要依赖代码执行环境、统计函数库和可视化组件完成落地；
在系统实现层面，模型又需要嵌入由状态对象、服务节点和消息通道组成的执行流程之中。
因此，工具调用不是简单的附加能力，而是把大语言模型从文本生成器转化为任务执行中枢的重要桥梁。

总体来看，Transformer架构、指令对齐、上下文学习、链式推理和工具增强共同构成了大语言模型在复杂数据任务中的基础能力体系。
这些研究为后续的查询生成、分析规划和系统化实现提供了统一的模型视角。

\section{面向可靠查询生成的关键技术}

Text-to-SQL旨在把用户自然语言问题转换为可在关系数据库上执行的SQL语句\cite{liu2025survey,kim2024text2sql_survey}。
若仅从任务定义看，它似乎只是“文本到代码”的转换问题；但在实际应用中，真正的难点在于自然语言、业务语义和物理数据库结构之间往往并不存在一一对应关系。
因此，可靠查询生成不仅要求模型“生成语句”，更要求系统具备识别业务术语、引入口径知识、组织样例上下文、借助执行反馈修复错误并从多候选中选择最优结果的能力\cite{王昊奋2024}。

\subsection{任务定义与核心挑战}

Text-to-SQL的基本输入是自然语言问题$Q$与目标数据库模式$S$，输出是一条语义等价、语法正确且可执行的SQL语句$y$\cite{liu2025survey}。
在传统公开基准中，研究重点主要集中于问题与Schema之间的语义对齐、复杂SQL结构生成以及跨库泛化能力。
随着大语言模型的引入，自然语言查询数据库技术逐渐从单纯的语句翻译扩展为“检索增强、结构约束、候选选择和执行验证”相结合的整体方案\cite{王昊奋2024}。

在垂直领域场景中，这一任务通常面临三类典型挑战。
第一，业务术语与物理字段之间可能存在一对多、多对一或缩写映射关系，导致模型即使理解了问题含义，也未必能准确定位字段。
第二，诸如同比增长率、预算执行率、结构占比等关键指标往往并非原始字段，而是依赖隐含计算公式和统计口径，若缺乏显式知识支撑，模型容易生成语法正确但业务错误的SQL。
第三，真实业务需求兼具高频固定模式与低频复杂长尾结构，单一路径生成往往难以同时兼顾效率与准确率。
沈宇等人\cite{沈宇2025}在交通系统智能问数研究中也表明，动态Schema对齐、领域知识注入和提示工程对于垂直领域NL2SQL落地具有显著作用。

已有研究从不同角度回应这些问题。
Seq2SQL\cite{zhong2017seq2sql}、SyntaxSQLNet\cite{yu2018syntaxsqlnet}开启了端到端神经生成路径，RAT-SQL\cite{wang2020ratsql}、BRIDGE\cite{lin2020bridge}、PICARD\cite{scholak2021picard}和RESDSQL\cite{li2023resdsql}则分别从关系建模、内容锚定、语法约束和Schema筛选等方面提升了查询生成可靠性。
在大模型阶段，DAIL-SQL\cite{gao2024text2sql}、DIN-SQL\cite{pourreza2023dinsql}、CHASE-SQL\cite{pourreza2024chase}、XiYan-SQL\cite{gao2024xiyan}与OpenSearch-SQL\cite{xie2025opensearch}说明，提示组织、样例检索、多路径生成与候选判别已经成为提升性能的重要手段。

\subsection{Schema/值检索与领域知识增强}

在可靠查询生成链路中，模式链接（Schema Linking）承担着“把自然语言中的实体、指标和条件映射到数据库结构”的基础作用。
RAT-SQL\cite{wang2020ratsql}通过关系感知编码显式建模问题Token与Schema元素之间的关系，BRIDGE\cite{lin2020bridge}利用数据库内容进行语义锚定，RESDSQL\cite{li2023resdsql}则通过排序与筛选把庞大Schema压缩为更适合生成的候选子模式。
这些工作共同揭示了一个事实：高质量SQL生成的前提不是“让模型直接猜”，而是先把与当前问题最相关的结构信息准确组织出来。

对于垂直领域系统，仅依赖原始表名、列名和注释通常仍然不够。
数据库中常见的字段缩写、别名表达、口径术语和隐含逻辑，往往无法被数据库模式显式刻画。
因此，结构化领域知识增强成为连接业务语言与物理数据库的必要桥梁。
一类方法通过补充别名、样例值、字段说明等语义提示提升模式可读性；
另一类方法则进一步把业务术语映射、逻辑公式和值域约束组织为可检索知识单元，并在在线推理时注入问题相关片段\cite{王昊奋2024,沈宇2025}。
这类思想本质上都是通过显式知识组织降低模型在字段绑定和业务理解上的歧义。

值检索与动态样例也是这一链路的重要组成部分。
对包含行政区划、项目名称、机构名称等实体的问题而言，仅靠模型记忆很难稳定命中正确取值，向量检索和近邻召回可以作为值匹配与样例选择的关键外部支撑。
在大模型Text-to-SQL中，DAIL-SQL\cite{gao2024text2sql}与OpenSearch-SQL\cite{xie2025opensearch}均证明了高相关动态样例对生成质量的重要影响。
这说明可靠查询生成并不是单轮提示行为，而是“检索 - 组织 - 生成”的联合过程。

\subsection{候选生成、执行反馈与结果选择}

大模型时代的Text-to-SQL系统逐渐从“单答案生成”转向“多候选生成 + 执行反馈 + 最优选择”。
其根本原因在于复杂SQL问题往往存在较高的推理不确定性：单次生成容易受到提示表述、随机采样和局部推理路径的影响，从而遗漏某个条件、错误绑定某个字段或采用错误聚合方式。
因此，多路径候选生成不再只是简单的重复采样，而是通过不同提示范式、不同分解路径和不同温度设置主动扩大候选空间。

CHASE-SQL\cite{pourreza2024chase}和XiYan-SQL\cite{gao2024xiyan}表明，多生成器和多推理路径能够显著提升复杂问题覆盖度；
Alpha-SQL\cite{li2025alpha}则进一步将SQL生成建模为搜索过程，说明复杂查询更适合采用“生成与评估交替”的策略。
另一方面，PICARD\cite{scholak2021picard}等工作说明，语法约束只能保证“写得像SQL”，却不能保证“问得对业务”，因此执行反馈成为进一步验证候选质量的重要机制。
当SQL真实执行后，系统可以根据报错信息、空结果、字段不匹配等反馈开展修复，从而把数据库引擎转化为事实校验器。

候选选择是保证最终输出质量的最后一道关口。
一致性投票适合处理低复杂度任务，但在存在业务逻辑冲突或执行表现差异时，仍需要更细粒度的判别机制。
近年来的偏好优化和候选判别方法表明，结合执行结果、语义一致性和任务约束进行成对比较，通常比简单投票更可靠\cite{pourreza2024chase,gao2024xiyan}。
总体来看，可靠查询生成正在从“单模型直接生成”转向“检索增强、知识注入、候选比较和执行验证”协同作用的综合范式\cite{王昊奋2024}。

\section{面向数据分析的规划与结果生成技术}

如果说Text-to-SQL解决的是“把数据取对”，那么面向数据分析的智能技术关注的则是“如何围绕已获取数据开展合适的分析并形成可理解的结果”。
与查询任务相比，数据分析任务具有更强的开放性：用户通常不会直接给出完整的分析步骤，而是以自然语言提出一个分析目标，例如“分析近两年各区财政收入变化并说明波动原因”。
系统因此需要先把需求转化为可执行规划，再调用查询、计算与可视化工具，最终输出图表、表格和文字洞察\cite{hong2024datainterpreter,gu2024analysts}。
随着大语言模型在政务与行业场景中的应用深化，数据分析与决策支持也逐渐成为重要落地方向\cite{王昀2024,王薇2025}。

\subsection{分析任务链条与中间规划}

已有研究普遍将自动数据分析过程拆分为需求理解、分析规划、工具执行和结果生成等阶段\cite{ma2023insightpilot,hong2024datainterpreter}。
其中，分析规划处于承上启下的位置：它既要把自然语言中的分析对象、指标、维度、时间范围和展示要求显式化，又要决定哪些步骤应交由查询模块完成，哪些步骤保留到分析层处理。
若缺乏这一中间层，系统就容易在复杂需求下出现“先取什么数据并不清楚”“展示和计算混在一起”“结果虽能生成但难以验证”等问题。

InsightPilot\cite{ma2023insightpilot}将数据探索拆解为一系列分析操作，说明分析任务本质上并非单步问答，而是带有过程结构的决策链。
Analysts\cite{gu2024analysts}则从更广泛视角评估了大语言模型作为数据分析者的可行性，指出模型在分析规划、统计执行和结果解释上的能力并不均衡，尤其需要外部约束来提高规划与输出质量。
在政务和行业场景中，这种约束更加重要，因为分析任务通常要遵循固定口径、固定展示习惯和更强的可复核要求\cite{王昀2024,王薇2025}。

基于此，结构化中间表示成为提高可控性的常见路径。
它的核心思想不是让模型一步产出最终结论，而是要求模型先输出一个受约束的分析计划，用以描述分析类型、关注指标、分析维度、筛选条件、计算方式以及展示规格。
常见的中间表示包括计划槽位、任务图、结构化脚本和可解析JSON对象等。
这类方式有助于把开放式需求转化为可校验、可复用、可执行的分析对象，从而提升系统整体稳定性。

\subsection{代码生成、分析算子与可视化组织}

在分析规划确定之后，系统还需要完成计算和表达两个层面的执行。
代码生成模型如Codex\cite{chen2021codex}、Code Llama\cite{roziere2023codellama}、StarCoder\cite{li2023starcoder}和WizardCoder\cite{luo2023wizardcoder}证明了大模型可以根据自然语言描述生成较高质量程序代码，这为自动统计分析、数据变换和图表绘制提供了技术前提。
但在实际系统中，直接让模型自由生成整段分析代码，往往会带来鲁棒性不足、重复逻辑过多和结果难以复核等问题。

因此，越来越多的系统开始采用“规划 + 算子 + 工具”的组织方式。
Data Interpreter\cite{hong2024datainterpreter}通过代码执行环境完成端到端分析，TaskWeaver\cite{qin2024taskweaver}强调代码优先与任务分工，LIDA\cite{dibia2023lida}和Chat2VIS\cite{maddigan2023chat2vis}则从可视化生成角度展示了模型如何驱动图表代码与视觉表达。
ChartGPT\cite{tian2024chartgpt}进一步说明，自然语言到图表之间并非简单模板映射，而是涉及图表类型选择、坐标映射和表达目标对齐。
与此同时，结构化输出评测研究表明，若缺少明确槽位和格式约束，模型在表格与结构数据生成上仍容易出现字段遗漏和格式漂移\cite{tang2024struc}。

这说明，数据分析结果生成不宜完全依赖自由代码生成，而更适合在分析算子和可视化规则层面引入明确边界。
在这种思路下，模型负责生成分析计划和调用意图，系统则依据预定义的轻量算子库完成排序、占比、透视、标注和图表组织等稳定操作。
这样既保留了大模型的灵活性，又提高了统计过程和展示结果的一致性。

\subsection{查询与分析分层协同}

分析系统中一个常见设计误区，是让分析模块重新承担数据库理解与SQL生成职责。
这样虽然看似减少了模块数量，但实际上会造成查询逻辑与分析逻辑混杂，导致一旦上游取数有误，后续图表和洞察都会建立在错误事实之上。
因此，越来越多的研究强调把“数据获取”和“数据分析”视作两个职责清晰但密切协同的阶段\cite{gu2024analysts,hong2024datainterpreter}。

分层协同的基本原则是：凡涉及数据库模式理解、复杂字段绑定、跨表关联和统计口径计算的操作，应尽量下沉到查询层完成；凡只依赖已获取结果的重组、排序、对比、占比和注释等操作，则由分析层负责。
这种划分的直接好处有三点。
第一，它复用了可靠查询模块在模式链接、逻辑增强和执行修复方面的已有能力，避免分析阶段重复犯错。
第二，它使分析规划能够围绕统一的中间表示展开，便于对查询约束和分析约束分别建模。
第三，它提升了系统整体可解释性，因为每一步究竟是“为了取数”还是“为了分析”都能够被明确追踪。

总体来看，面向数据分析的智能技术正逐步从“直接生成结论”转向“结构化规划、工具执行和结果组织”相结合的实现路径。
这种分层化、可执行、可复核的分析范式，已经成为复杂数据应用的重要发展方向\cite{王昀2024,王薇2025}。

\section{面向系统实现的工作流编排与服务协同}

当查询与分析能力从单点方法发展为可交互系统后，研究关注点便进一步延伸到工程实现层面。
一个实用系统不仅要给出结果，还需要组织任务流程、维护状态、回传异常并沉淀中间产物，以支持复核、追问和问题定位。
这使得工作流编排、状态管理、服务协同和角色分工成为大语言模型应用落地时不可回避的关键技术\cite{孟天广2021,张聪丛2024,王薇2025}。

\subsection{工作流图编排、状态管理与异常回传}

大语言模型应用逐步走向复杂系统后，单轮提示已经难以承载完整业务流程。
一个典型任务往往需要经历需求解析、知识检索、SQL生成、执行验证、结果封装和前端回传等多个阶段。
因此，工作流图编排成为组织模型能力与系统组件的重要方式。
从原理上看，工作流图把任务拆分为若干节点，每个节点负责相对单一的职责，节点之间通过显式状态对象和条件边连接，从而形成可追踪、可恢复的执行路径。

ReAct\cite{yao2023react}提供了“推理 - 行动 - 观察”的通用闭环思想，表明模型驱动任务执行更适合放入可循环、可反馈的流程而非一次性输出。
对系统实现而言，这一思想需要进一步工程化：系统要把模型调用、数据库执行、文件生成和消息回传等步骤纳入统一状态机中，并对成功、失败、重试和终止等状态进行管理。
周祥生等人\cite{周祥生2025}关于大模型推理任务调度的研究也说明，随着模型能力嵌入真实业务链路，任务优先级、资源分配和状态调度会直接影响系统稳定性与用户体验。
同时，政务垂类大模型的应用实践表明，只有把模型能力放入可管理、可审计的服务流程中，才能兼顾灵活性与可控性\cite{吴重庆2025,兰洪浩2025}。

\subsection{轻量级Agent协同与角色分工}

在复杂任务中，不同节点往往承担不同角色：有的负责路由请求，有的负责生成查询，有的负责组织分析，有的负责汇总结果。
这类“角色分工 + 消息协同”的设计与近年来的Agent研究存在天然联系。
AutoGen\cite{wu2023autogen}强调通过多角色对话组织工具调用与任务分工，MetaGPT\cite{hong2024metagpt}强调基于标准操作流程的结构化协同，CAMEL\cite{li2023camel}则展示了角色扮演框架下的任务协商能力。
郭先会等人\cite{郭先会2024}从中文综述角度指出，大语言模型智能体通常围绕任务规划、工具执行、记忆组织和评测机制展开，其工程价值很大程度上取决于角色职责是否清晰、工具接口是否稳定。

不过，在真实业务场景中，更具可操作性的往往不是“重型多智能体社会模拟”，而是轻量级服务协同。
也就是说，系统不必追求大量Agent之间的开放式对话，而是更强调把查询、分析、调度和结果组织等职责稳定封装为可复用服务角色。
李轶夫\cite{李轶夫2025}指出，生成式大模型智能体真正形成工程价值的前提，是把角色能力、工具接口和业务场景有机结合，而不是停留在概念展示层面。
相关政务实践也说明，分层设计、模块化集成和模型服务化部署，是保障系统稳定落地的重要路径\cite{王薇2025,吴重庆2025,兰洪浩2025}。

综合来看，工作流编排负责组织任务过程，状态管理负责约束执行边界，轻量级Agent协同则有助于明确角色职责并降低上下文混杂。
这些研究共同构成了大语言模型系统从“能回答”走向“能稳定运行”的工程技术基础。

\section{本章小结}

本章围绕与本文研究相关的理论与技术，对大语言模型基础、可靠查询生成、数据分析规划与结果生成，以及系统编排与服务协同等内容进行了综述。

在模型基础方面，Transformer架构、指令对齐、上下文学习、链式推理和工具增强共同塑造了大语言模型处理复杂数据任务的核心能力。
在查询生成方面，Schema链接、领域知识增强、值检索、动态样例、多候选生成和执行反馈构成了提高NL2SQL可靠性的关键技术链。
在数据分析方面，结构化规划、代码执行、分析算子和查询分析分层协同体现了从开放式需求到规范化分析结果的主流实现思路。
在系统实现方面，工作流编排、状态管理与轻量级Agent协同则为复杂智能系统的稳定运行、可追溯性和工程落地提供了重要支撑。

总体而言，上述研究共同描绘了面向复杂数据应用的技术演进方向，即从单一模型生成逐步走向“知识增强、工具执行、流程组织和结果验证”协同作用的综合范式。
